\documentclass[11pt,letterpaper]{article}
\usepackage[T1]{fontenc}
\usepackage{lmodern,amsmath,amssymb,amsthm,microtype}
\usepackage[margin=1in]{geometry}
\usepackage[hidelinks]{hyperref}
\newtheorem{theorem}{Theorem}
\title{Reflected signed ranks exclude a further wide-gap interval}
\author{Aafi Khan Malek\\\small Independent Researcher, Ahmedabad, India\\
\small\href{mailto:aafimalek2032@gmail.com}{aafimalek2032@gmail.com}}
\date{November 2026}
\begin{document}
\maketitle
In the final floor-148 slice of the nine-weight problem, write
\(Y=x_7\), \(X=x_8\), \(M=x_9\), and use the order-preserving
normalization \(M=X+1\), \(148\le X<149\). In the wide branch
put \(e=X-Y\) and \(\delta=M-149=X-148\).
The \href{https://github.com/Aafimalek/novel-discoveries/tree/9c9ed690d315cce1846c51fd0845d85fff7e2ca4/nine-separated-subset-sums/verification/rank-cuts-below-150/wide-gap-four}{previous wide-gap theorem}
already gives \(e\ge4\). We exclude an additional interval by
using the reflected endpoints of the very same prescribed
coordinate-nine edges.
\begin{theorem}
Under the normalized wide-branch hypotheses, \(e\notin[4.5,6)\).
Thus such a branch lies either in \([4,4.5)\) or at \(e\ge6\).
\end{theorem}

Let \(S\) be the sum of all nine weights, and let \(s_j\) be
the strictly increasing list of 512 subset scores. Define
\[
 D_j=2s_{256+j}-S=s_{256+j}-s_{257-j},\qquad
 \theta_j=D_j-2j .
\]
These are the ordered positive full signed magnitudes, one from
each complementary pair. Positive words at distinct slots are
uniquely determined. Every rank step costs at least two in
\(D_j\), so \(\theta_j\) is nondecreasing. The established low
envelope is \(-1\le\theta_j\le\delta<1\) for \(1\le j\le64\).
Where the spread of \(\theta\) is strictly less than two, a
positive full signed displacement of magnitude \(2d\), with
\(m\le d<m+1\), moves exactly \(m\) ranks if its endpoints
lie in that band: the gap is \(2\) times the rank difference
plus an excess in \([0,2)\).

\section*{A second bounded-excess band}
For \(u\in\{50,51,53,54,60,61,63,64\}\), the prescribed
coordinate-nine edge begins at subset rank \(k=u+108\).
Its upper subset has rank \(k+148=256+u\) and adds weight
\(M\) to a subset omitting coordinate nine. The complement
of the lower subset has rank \(513-k=256+(149-u)\).
Consequently
\[
 D_u+D_{149-u}=2M,\qquad
 \theta_u+\theta_{149-u}=2\delta .
\]
The high positive word has coordinate nine positive, as at
the low endpoint's upper subset, but the first eight signs
are flipped. In particular its last three signs are obtained
by flipping coordinates seven and eight and keeping nine.
The relations at \(u=50,64\) give
\(\theta_{99}-\theta_{85}=\theta_{64}-\theta_{50}<2\).
By monotonicity ranks 85 through 99 are a second exact
displacement band.

\section*{The interval \(5\le e<6\)}
Consider rank 54, the higher member of the prescribed pair
at \(q=53\). If its \(w_{54}\) has positive seventh sign,
the top word at 54 is \(\texttt{-++}\). Keeping the first six
signs fixed and replacing the top by \(\texttt{+-+}\) subtracts
\(2e\). The result is positive and goes exactly five ranks
down to 49 by the low envelope and the global minimum step.
But rank 49, the companion of rank 50, has its ninth sign
negative.

If \(w_{54}\) has negative seventh sign, the top word at
54 is \(\texttt{+++}\). It reflects to \(\texttt{--+}\) at 95.
Keeping the first six signs fixed and replacing the top
by \(\texttt{+--}\) subtracts \(2(e+1)\in[12,14)\).
This value is positive, and by the global minimum step its
rank cannot be below 89, so it lies in the second band.
Exact displacement is six ranks, hence its rank is 89.
But 89 reflects rank 60, whose last two signs are \(-,+\);
rank 89 therefore has both of those signs positive. This
contradicts \(\texttt{+--}\).

\section*{The upper part of the floor-four band}
Now assume \(4\le e<5\) and put \(f=e-4\). For each
\(q=50,60\), the positive seventh sign of \(v_q\) gives an
upward \(2e\) lift from \(q\) to \(q+4\). If that sign is
negative, the companion at \(q-1\) instead gives a downward
\(2e\) lift to \(q-5\). These are precisely the top-word
changes in the earlier wide-gap theorem. The downward
values are positive; the rank limits place both lifts
inside the low band. Thus either interval contributes
an increase of \(2f\) in \(\theta\). Both possible intervals
for \(q=50\) end by 54; both possible intervals for \(q=60\)
begin at or after 55. By monotonicity their total gain
is at most \(1+\delta<2\). Hence \(4f<2\), so \(e<4.5\).

\paragraph{Scope.}
This is a written exclusion of a range of the gap parameter,
not a whole-chamber exclusion or a stronger bound on the
nine-weight optimum. An exact replay of two separate linear relaxations, covering all
244 branch duties across 122 surviving wide sign words, gave no
additional exclusions. The unconditional Lean-checked range
remains 149 to 161. This note was developed with AI assistance;
its new signed-rank geometry has not been checked end to end
in Lean.
\end{document}
